\section{Resolución de los problemas más relevantes}

En esta sección se describen los problemas más relevantes encontrados durante el desarrollo y cómo se han resuelto.

\subsection{Conversión entre \texttt{cv::Mat} y \texttt{QImage}}

Problema: OpenCV usa matrices en orden BGR/BGRA y Qt espera RGB/RGBA en ciertos formatos; además la memoria subyacente puede ser gestionada por OpenCV, lo que puede provocar accesos a memoria liberada si no se copian los datos.

Solución: Antes de crear un \texttt{QImage} se convierte el orden de canales (p. ej. \texttt{COLOR\_BGR2RGB} o \texttt{COLOR\_BGRA2RGBA}) y se realiza una copia con \texttt{QImage::copy()} cuando se va a transferir al portapapeles o a Qt para garantizar la propiedad de los datos.

\subsection{Gestión de múltiples ventanas}

Problema: controlar el foco y el orden de múltiples ventanas creadas con HighGUI de OpenCV.

Solución: se utiliza el array global \texttt{foto[]} con un campo \texttt{orden} que se incrementa cada vez que una ventana obtiene foco; la función \texttt{foto\_activa()} devuelve la ventana con mayor orden.

\subsection{Captura de cámara y sincronización}

Problema: capturar frames desde la cámara y permitir operaciones como acumulado de media sin bloquear la interfaz.

Solución: el módulo \texttt{video.cpp} encapsula la lógica de captura con \texttt{VideoCapture}. Las funciones devuelven controles sencillos (por ejemplo, número de ms entre frames) y se usan ventanas de OpenCV para previsualizar; las operaciones de procesado (media, rotación, exportación) se realizan en funciones separadas para evitar entrelazar la lógica de captura y la de procesamiento.

\subsection{Compatibilidad de formatos al copiar al portapapeles}

Problema: el portapapeles del sistema puede requerir diferentes formatos según la aplicación destino.

Solución: se soportan conversiones para imágenes de 1, 3 y 4 canales y se colocan en el portapapeles como \texttt{QImage}; en Windows esto permite pegar en Paint y otras aplicaciones que aceptan imágenes.
